\section{Discussion (Обсуждение)}
\label{sec:discussion}

\subsection{Концепция Zero Alignment Tax: преодоление компромисса}
В литературе по выравниванию языковых моделей (RLHF) и рекомендательных систем широко задокументирован так называемый «налог на выравнивание» (\textit{Alignment Tax})~\cite{ouyang2022training}: попытка оптимизировать внешнее вознаграждение неизбежно приводит к деградации базовых способностей модели к предсказанию распределения данных.

В наших экспериментах на Yambda-50M чистый DPO без регуляризации ($\gamma=0.0$, Trial \#1 в Таблице~\ref{tab:sweep_results}) наглядно подтвердил эту опасность: модель быстро максимизирует Win Rate предпочтений до 67.19\% и раздувает скрытую маржу до $+0.1735$, однако валидационный HR@10 катастрофически обрушивается с 0.0787 до 0.0688 (относительное падение на $-12.5\%$, а на полном тестовом сете HR@10 падает до $7.01\%$). Модель начинает эксплуатировать глобальные паттерны длинных треков в ущерб семантической релевантности конкретного пользователя.

Введение вспомогательного SFT-якоря ($\gamma=0.10$, Уравнение~\ref{eq:dpo_sft}) устраняет этот компромисс. Модель удерживается в области правдоподобных языковых продолжений, что обеспечивает \textbf{Zero Alignment Tax}: точность ранжирования следующего трека не только не падает, но и достоверно возрастает (+0.67\% HR@10, $p < 0.0001$).

\subsection{Анатомия эффекта «полупроцента»: практическая vs статистическая значимость}
При интерпретации результатов индустриальных рекомендательных систем критически важно разделять статистическую неслучайность и величину прикладного эффекта:
\begin{enumerate}
    \item \textbf{Практический масштаб:} Рост HR@10 с 19.23\% до 19.89\% (+0.67\% в абсолюте, +3.48\% в относителе) означает, что на отложенной тестовой выборке из 9\,048 пользователей базовая модель успешно нашла следующий трек в Top-10 для 1\,740 человек, тогда как DPO нашла его для 1\,800 человек. Чистый выигрыш составил ровно \textbf{+60 пользователей}.
    \item \textbf{Анализ дискордантных пар МакНемара:} DPO исправил ошибку базовой модели на 125 пользователях ($n_{01}$), однако сам ошибся там, где Base угадал, на 65 пользователях ($n_{10}$). Итоговый размер эффекта Коэна невелик ($h \approx 0.04$).
    \item \textbf{Статистическая неслучайность:} Все 5\,000 бутстрэп-итераций дают строго положительную дельту (95\% ДИ $[+0.0038, +0.0097]$, $p < 0.0001$), а тест МакНемара дает $\chi^2 = 18.32, p = 1.9 \times 10^{-5}$.
\end{enumerate}

Мы открыто констатируем: метрика Next-Item HR@K — это суррогатная метрика задачи MLE, а не истинная цель обучения с подкреплением. Главная победа DPO выражается не в полупроцентном сдвиге угадывания токена, а в качественном отсечении скипов, росте непрерывной сессии (+1.5 с) и оздоровлении каталога ($-12.1\%$ зависимости от поп-хитов).

\subsection{Почему DPO превзошёл двухстадийные подходы (Critic и PPO)}
Эксперименты наглядно выявили причины неэффективности двухстадийных методов в оффлайн-логах:
\begin{itemize}
    \item \textbf{Информационное плато оффлайн-критика:} В датасете Yambda доля скипов составляет 23.8\%. Теоретическая энтропия Бернулли $H(p) \approx 0.5488$. Функция потерь критика застряла на отметке $0.523$, что означает неспособность критика предсказывать индивидуальную оценку пользователя точнее тривиальной априорной вероятности популяции.
    \item \textbf{Бутылочное горлышко замороженных эмбеддингов:} MLP-критик (509K параметров) обучался поверх замороженных эмбеддингов $wte$, оптимизированных под совместную встречаемость токенов. Он не имел доступа к акустическим свойствам музыки и не смог сформировать разделяющую гиперплоскость для 98.5K треков.
    \item \textbf{Reward Hacking в PPO:} При оптимизации актёра против несовершенного критика PPO быстро начинает эксплуатировать локальные артефакты функции ценности, рекомендуя треки с завышенным вознаграждением в ущерб смысловой связности сессии (HR@10 просел до 18.60\%).
    \item \textbf{Преимущество DPO:} DPO полностью обходится без явной функции ценности, сквозным образом оптимизируя все 43.2M параметров самого трансформера через устойчивую сигмоидальную маржу логарифмов отношений правдоподобия.
\end{itemize}

\subsection{Качественный ручной аудит пользовательских траекторий}
Ручной аудит тестовых примеров иллюстрирует качественное изменение поведения:
\begin{itemize}
    \item \textbf{Кейс 1 (User 337300, длинная сессия):} Пользователь слушает треки длительностью 3.5--4 минуты. Истинный таргет `4515849` прослушан целиком (240.0 с). Базовый GPT-2 выдал его на Rank \#1 с уверенностью 63.14\%. DPO усилил уверенность до \textbf{69.79\% (+6.65\%)} за счёт позитивной истории удержания.
    \item \textbf{Кейс 2 (User 366500, серия скипов):} Таргет `410531` был переключен пользователем за \textbf{0.0 секунд} (мгновенный отказ). Базовая модель классифицировала отказ от рекомендации как ошибку, тогда как DPO целенаправленно понизил ранг этого трека, предотвратив негативный пользовательский опыт.
\end{itemize}

\subsection{Ограничения исследования (Threats to Validity)}
\begin{enumerate}
    \item \textbf{Counterfactual Gap:} В оффлайн-логах невозможно измерить реальную реакцию пользователя на трек, которого не было в истории (не наблюдаемый контрфактический исход).
    \item \textbf{Отсутствие сырого аудио:} Модель оперирует категориальными идентификаторами треков (ID tokens) без учёта аудио-спектрограмм и темпо-ритмических признаков.
    \item \textbf{Статичность предпочтений:} Симуляция сессий не учитывает контекст времени суток, геолокации и типа устройства.
\end{enumerate}

\subsection{Дизайн промышленного онлайн A/B тестирования}
Для валидации разработанного подхода в промышленной среде рекомендуется следующий протокол A/B эксперимента:
\begin{itemize}
    \item \textbf{Сплит и длительность:} 50/50 случайное распределение пользователей по идентификаторам сессий на протяжении не менее 3--4 недель (для нивелирования эффекта новизны).
    \item \textbf{Первичная бизнес-метрика:} Среднее время непрерывного прослушивания за сессию (\textit{Session Playtime Duration}) и доля скипов в первые 30 секунд ($P_{\text{skip}}$).
    \item \textbf{Долгосрочные прокси:} Retention 7-го и 30-го дня (D7/D30), а также вероятность возврата в сервис в течение 24 часов после завершения текущей сессии ($Y_{\text{return}} \in \{0, 1\}$).
    \item \textbf{Guardrail-метрики:} Задержка генерации следующего трека ($p99 < 35$ мс) и общая частота обращений к поиску (как сигнал потери интереса к радиопотоку).
\end{itemize}
